\documentclass[11pt]{amsart}
\usepackage{iftex}
\ifPDFTeX\usepackage{cmap}\fi
\usepackage[T1]{fontenc}
\usepackage{lmodern}
\usepackage{amssymb,amsthm,array,booktabs}
\usepackage[hidelinks]{hyperref}
\hypersetup{
  pdftitle={A line-direction lemma for the transversal difference number of (Z/p\textasciicircum 2Z)\textasciicircum 2},
  pdfauthor={Sukhraj Sachdeva},
  pdfkeywords={transversal, difference set, sumset, finite abelian group, transversal difference number}
}

\newtheorem{theorem}{Theorem}[section]
\newtheorem{lemma}[theorem]{Lemma}
\newtheorem{proposition}[theorem]{Proposition}
\newtheorem{corollary}[theorem]{Corollary}
\newtheorem{conjecture}[theorem]{Conjecture}
\theoremstyle{definition}\newtheorem{definition}[theorem]{Definition}
\theoremstyle{remark}\newtheorem{remark}[theorem]{Remark}

\newcommand{\F}{\mathbb{F}}
\newcommand{\Z}{\mathbb{Z}}
\newcommand{\abs}[1]{\lvert #1\rvert}
\newcommand{\gen}[1]{\langle #1\rangle}
\newcommand{\file}[1]{\nolinkurl{#1}}
\DeclareMathOperator{\id}{id}
\raggedbottom

\title[A line-direction lemma for transversal differences]{A line-direction lemma for the transversal difference number of $(\mathbb{Z}/p^2\mathbb{Z})^2$}
\author{Sukhraj Sachdeva}
\date{2026-11-27. DOI: \href{https://doi.org/10.5281/zenodo.24594748}{10.5281/zenodo.24594748}}
\subjclass[2020]{Primary 11B13; Secondary 11P70, 20K01}
\keywords{transversal, difference set, sumset, finite abelian group, transversal difference number}

\begin{document}
\begin{abstract}
Let \(p\) be an odd prime, \(G=(\Z/p^2\Z)^2\) and \(H=pG\).
Barcau, Pa\c{s}ol and \c{T}urca\c{s} conjectured that every transversal \(T\) of \(H\) in \(G\) satisfies \(\abs{T-T}\ge(2p-1)^2\).
They write each transversal through a function \(f\colon\F_p^2\to\F_p^2\), and for each \(u\in\F_p^2\) the differences of \(T\) that reduce to \(u\) modulo \(p\) correspond to a set \(A_u(f)\subseteq\F_p^2\).
The conjectured bound holds whenever \(A_u(f)\) lies on an affine line for some \(u\ne0\).
This lemma was also obtained, earlier on the same day, by Leshner using the AI system Codex \cite{leshner}.
We give an independent proof, which avoids Kneser's theorem, and we claim no priority.
As consequences, every transversal has at least \(3p^2-2\) differences, the case \(p=3\) (known by computation) follows by hand, and the conjecture is equivalent to its restriction to functions for which no such \(u\) exists.
This remaining case is not settled here.
A computer-assisted proof of the case \(p=7\) has been claimed in \cite{leshner}; we have not checked it.
This human-directed research was produced with the help of AI tools and checked by AI referees; additional verification may still be required.
\end{abstract}
\maketitle

\section{Introduction}\label{sec:intro}

\subsection{The problem}
Let \(G\) be a finite abelian group and \(H\) a subgroup.
A \emph{transversal} of \(H\) in \(G\) is a set \(T\subseteq G\) that contains exactly one element of each coset of \(H\).
Its \emph{difference set} is \(T-T=\{t-t'\mid t,t'\in T\}\).
Barcau, Pa\c{s}ol and \c{T}urca\c{s} \cite{BPT} study the \emph{transversal difference number}
\[
\delta(G,H)=\min\{\abs{T-T}\mid T \text{ a transversal of } H \text{ in } G\}.
\]
For a prime \(p\) put \(G_p=(\Z/p^2\Z)^2\) and \(H_p=pG_p\).
So \(H_p\) is the subgroup of elements whose two coordinates are multiples of \(p\).
The classes of the integer points of the ``box'' \(\{0,\dots,p-1\}^2\) form a transversal of \(H_p\) in \(G_p\) with exactly \((2p-1)^2\) differences.
Hence \(\delta(G_p,H_p)\le(2p-1)^2\) \cite[Proposition~5.3]{BPT}.
The source conjectures that the box is optimal.

\begin{conjecture}[{\cite[Conjecture~6.2]{BPT}}]\label{conj:BPT}
For every odd prime \(p\),
\[
\delta\bigl((\Z/p^2\Z)^2,\,p(\Z/p^2\Z)^2\bigr)=(2p-1)^2 .
\]
\end{conjecture}

The source proves \(\delta(G_p,H_p)\ge3p^2-p-1\) for every odd prime \(p\) \cite[Theorem~6.5]{BPT}.
It reports a computer verification of the conjecture for \(p=3\) and \(p=5\) \cite[\S6.4]{BPT}, with code in \cite{BPTcode}.
It also proves that random transversals satisfy the conjectured bound with probability tending to \(1\), and that transversals given by a fixed pair of integer polynomials satisfy it for all large \(p\) \cite[Theorems~7.1 and~7.2]{BPT}.

\subsection{Results}
Section~\ref{sec:setting} recalls the coordinates used in the source.
In these coordinates, each transversal of \(H_p\) in \(G_p\) is a set \(T_f\) for exactly one function \(f\colon\F_p^2\to\F_p^2\).
For each \(u\in\F_p^2\), the source defines a set \(A_u(f)\subseteq\F_p^2\), see \eqref{eq:Au}.
Its elements correspond one-to-one to the differences in \(T_f-T_f\) that reduce to \(u\) modulo \(p\).
We call a nonzero \(u\) a \emph{line direction} of \(f\) if \(A_u(f)\) lies on an affine line of \(\F_p^2\).

The main statement, Theorem~\ref{thm:main}, says that one line direction forces \(\abs{T_f-T_f}\ge(2p-1)^2\).
This part was obtained earlier on the same day by Leshner with Codex \cite{leshner} (see Section~\ref{sec:prior}).
Theorem~\ref{thm:main} also gives a lower bound for the number of differences over each coset of the line \(\{\lambda u\mid\lambda\in\F_p\}\).
And it characterises line directions: \(f\) has one exactly when \(T_f\) is contained in a union of at most \(p\) cosets of a cyclic subgroup of order \(p^2\).

Three consequences follow (Section~\ref{sec:statements}).
\begin{itemize}
\item \(\delta(G_p,H_p)\ge3p^2-2\) for every odd prime \(p\) (Corollary~\ref{cor:lower}).
\item Conjecture~\ref{conj:BPT} at \(p\) is equivalent to its restriction to functions with no line direction (Corollary~\ref{cor:reduction}).
\item \(\delta(G_3,H_3)=25\), proved without computation (Corollary~\ref{cor:p3}).
\end{itemize}
We do not settle the functions with no line direction.
A computer-assisted proof of the case \(p=7\) has been claimed in \cite{leshner}; we have not checked it.
Beyond that, the conjecture remains open as far as we know.

\subsection{Prior work and credit}\label{sec:prior}
The statement of Theorem~\ref{thm:main}(a) was posted publicly earlier on the same day as this work, 2026-11-27, by Leshner with Codex \cite{leshner}, as a ``projection bound for every odd prime''.
Their proof changes basis, applies a shear, projects onto a cyclic transversal, and uses Kneser's addition theorem.
Our work was done independently and was posted later the same day.
Our proof covers \(T_f\) by cosets of a cyclic subgroup and uses an elementary lemma on two transversals (Lemma~\ref{lem:cross}) in place of Kneser's theorem.
Their lemma, together with the count in the proof of Corollary~\ref{cor:lower}, also gives Corollaries~\ref{cor:lower}--\ref{cor:p3}.
The same pull request additionally claims a computer-assisted proof that \(\delta(G_7,H_7)=169\); we have not checked that claim.
We claim no priority for Theorem~\ref{thm:main}(a) or for any of the corollaries.

In the source, two linearly independent directions \(u,v\) with \(\abs{A_u(f)}=\abs{A_v(f)}=2\) force \(\abs{T_f-T_f}\ge(2p-1)^2\) \cite[Lemma~6.4]{BPT}.
The value \(\delta(G_3,H_3)=25\) is already in the source, by computer \cite[\S6.4]{BPT}.

Unpublished working notes in the public repository chy4pro/automath \cite{automath} also concern this conjecture.
They consist of a numerical log, an unreviewed proof sketch, and a status note.
The status note calls the following statement a ``provisional theorem'': a single direction \(u\ne0\) with \(\abs{A_u(f)}=2\) forces \(\abs{T_f-T_f}\ge(2p-1)^2\).
The notes state this case together with its consequences.
These are the column bound \(4p-2\) of Theorem~\ref{thm:main}(b) in that case, a count giving \(3p^2-2\) for the remaining functions and hence the case \(p=3\) by hand, and a reduction of the conjecture to the functions with \(\abs{A_w(f)}\ge3\) for every \(w\ne0\).
The log records that remaining class as open.
We did not locate a complete written proof in the notes, and we have not compared our proof with their sketch.
A two-element set always lies on a line.
So the statement of the notes is the special case \(\abs{A_u(f)}=2\) of Theorem~\ref{thm:main}(a).
Theorem~\ref{thm:main} allows a collinear set \(A_u(f)\) of any size and adds the covering form.
Its residual class in Corollary~\ref{cor:reduction} is therefore smaller.

An issue opened on 2026-09-17 in the public repository trureturing of the-omega-institute \cite{trureturing} isolated the remaining inequality in the case where the directions \(u\) with \(\abs{A_u(f)}=2\) lie on one line through \(0\).

The setting, the conjecture, the fibre formula and the cycle identity used here are due to Barcau, Pa\c{s}ol and \c{T}urca\c{s} \cite{BPT}, as are the two-direction lemma and the bound \(3p^2-p-1\).
We make no claim of priority or novelty.

Sections~\ref{sec:setting}--\ref{sec:corollaries} prove every statement by hand, with no computational step; Section~\ref{sec:computer} describes computer checks, which are evidence only, and the state of formalization.

\section{Setting and notation}\label{sec:setting}

Throughout, \(p\) is a prime and \(\F_p=\Z/p\Z\).
In all statements \(p\) is odd, as in the source.
The proofs never use that \(p\) is odd (Remark~\ref{rem:p2}).
We write \(G=G_p=(\Z/p^2\Z)^2\), a group of order \(p^4\), and \(H=H_p=pG\), a subgroup of order \(p^2\).
For an element \(g\) of an abelian group and an integer \(j\), \(jg\) is the \(j\)-fold sum.
We write \(\gen{g}\) for the cyclic subgroup generated by \(g\); the \emph{order} of \(g\) is the number of elements of \(\gen{g}\).

Let \(\pi\colon G\to\F_p^2\) be coordinatewise reduction modulo \(p\).
It is a surjective homomorphism, and its kernel is \(H\): a pair of residues modulo \(p^2\) reduces to \(0\) modulo \(p\) exactly when both are multiples of \(p\).
So a set \(T\subseteq G\) is a transversal of \(H\) if and only if the restriction of \(\pi\) to \(T\) is a bijection onto \(\F_p^2\).
In particular, a transversal has \(p^2\) elements.

For \(x\in\F_p^2\), its \emph{standard lift} \([x]\in\{0,1,\dots,p-1\}^2\subseteq\Z^2\) is the vector of least non-negative representatives.
The same symbol denotes the class of \([x]\) in \(G\).
Define \(\iota\colon\F_p^2\to G\) by \(\iota(a)=pm\bmod p^2\), where \(m\in\Z^2\) is any integer vector with \(m\equiv a\pmod p\).
For a function \(f\colon\F_p^2\to\F_p^2\), put
\[
t_f(x)=[x]+\iota(f(x))\in G,\qquad T_f=\{t_f(x)\mid x\in\F_p^2\}.
\]
We write \(t\) for \(t_f\) when \(f\) is fixed.
For \(u,x\in\F_p^2\) the source defines the \emph{carry} \(c_u(x)\), the \emph{corrected difference} \(d_uf(x)\) and its \emph{image} \(A_u(f)\):
\begin{align}
c_u(x)&=\frac{[x+u]-[x]-[u]}{p}\in\F_p^2,\label{eq:carry}\\
d_uf(x)&=f(x+u)-f(x)+c_u(x),\label{eq:du}\\
A_u(f)&=\{d_uf(x)\mid x\in\F_p^2\}.\label{eq:Au}
\end{align}
The numerator in \eqref{eq:carry} is computed in \(\Z^2\).
Each of its coordinates is \(0\) or \(-p\).
So each coordinate of \(c_u(x)\) is \(0\) or \(-1\), read modulo \(p\).

For \(s\in\F_p^2\) let \(\gen{s}=\{\lambda s\mid\lambda\in\F_p\}\).
An \emph{affine line} is a set \(a+\gen{s}\) with \(a\in\F_p^2\) and \(s\ne0\); it has \(p\) points.
For \(u\ne0\), a \emph{\(u\)-line} is a coset \(x+\gen{u}\), that is, an affine line parallel to \(u\).
Any set of at most two points lies on an affine line.
Indeed, one point \(a\) lies on \(a+\gen{s}\) for any \(s\ne0\), and two points \(a\ne b\) lie on \(a+\gen{b-a}\).

\begin{definition}\label{def:line}
A \emph{line direction} of \(f\) is a vector \(u\in\F_p^2\setminus\{0\}\) such that \(A_u(f)\) is contained in an affine line.
\end{definition}

For \(u\ne0\) and \(v\in\F_p^2\), the \emph{column sum} of the \(u\)-line \(v+\gen{u}\) is the total size of the sets \(A_w(f)\) over its points \(w\), that is, \(\sum_{k\in\F_p}\abs{A_{ku+v}(f)}\).

The next lemma collects three facts about \(\iota\) and orders.
\begin{lemma}\label{lem:facts}
\begin{enumerate}
\item[(a)] The map \(\iota\) is well defined, additive and injective, and \(\iota(\F_p^2)=H\). In particular \(\iota(a)=p[a]\).
\item[(b)] \(pg=\iota(\pi(g))\) for every \(g\in G\). Hence \(pg=0\) if and only if \(g\in H\).
\item[(c)] Every \(g\in G\setminus H\) has order \(p^2\).
\end{enumerate}
\end{lemma}

\begin{proof}
(a) If \(m\equiv m'\pmod p\), then \(pm\equiv pm'\pmod{p^2}\), so \(\iota\) is well defined.
If \(m,m'\) lift \(a,a'\), then \(m+m'\) lifts \(a+a'\), so \(\iota\) is additive.
If \(pm\equiv0\pmod{p^2}\), then \(m\equiv0\pmod p\), so \(\iota\) is injective.
The image is \(\{pm\bmod p^2\mid m\in\Z^2\}=pG=H\).
Since \([a]\) lifts \(a\), we have \(\iota(a)=p[a]\).

(b) Write \(g=m\bmod p^2\) with \(m\in\Z^2\).
Then \(\pi(g)=m\bmod p\) and \(pg=pm\bmod p^2=\iota(\pi(g))\).
By (a), \(\iota(\pi(g))=0\) exactly when \(\pi(g)=0\), that is, when \(g\in H\).

(c) Since \(p^2g=0\), the order of \(g\) divides \(p^2\).
By (b), \(pg\ne0\), so the order does not divide \(p\).
\end{proof}

\begin{lemma}\label{lem:Tf}
The transversals of \(H\) in \(G\) are exactly the sets \(T_f\), and each transversal is \(T_f\) for exactly one function \(f\).
\end{lemma}

\begin{proof}
Since \(\pi(t_f(x))=x\) for every \(x\), the set \(T_f\) is a transversal.
Conversely, let \(T\) be a transversal and \(x\in\F_p^2\).
The unique element of \(T\) with image \(x\) under \(\pi\) is \([x]+h\) for a unique \(h\in H\).
By Lemma~\ref{lem:facts}(a), \(h=\iota(f(x))\) for a unique \(f(x)\in\F_p^2\).
\end{proof}

The source also normalises \(f(0)=0\), that is, \(0\in T\).
Nothing below needs this.

\section{Statements}\label{sec:statements}

\begin{theorem}[line-direction lemma]\label{thm:main}
Let \(p\) be an odd prime, \(f\colon\F_p^2\to\F_p^2\) any function and \(T=T_f\).
In parts (a)--(c), suppose that \(u\ne0\) is a line direction of \(f\).
\begin{enumerate}
\item[(a)] \(\displaystyle\abs{T-T}=\sum_{w\in\F_p^2}\abs{A_w(f)}\ge(2p-1)^2\).
\item[(b)] For every \(v\in\F_p^2\setminus\gen{u}\), the column sum of \(v+\gen{u}\) satisfies
\[
\sum_{k\in\F_p}\abs{A_{ku+v}(f)}\ge4p-2 .
\]
Moreover, \(\sum_{k\in\F_p}\abs{A_{ku}(f)}\ge2p-1\).
\item[(c)] \(T\) is contained in a union of \(p\) cosets of a cyclic subgroup \(C\le G\) of order \(p^2\) with \(\pi(C)=\gen{u}\).
\item[(d)] Conversely, suppose that \(T\) is contained in a union of at most \(p\) cosets of a cyclic subgroup \(C\le G\) of order \(p^2\).
Then every nonzero \(w\in\pi(C)\) is a line direction of \(f\).
\item[(e)] Hence \(f\) has a line direction if and only if \(T\) is contained in a union of at most \(p\) cosets of some cyclic subgroup of \(G\) of order \(p^2\).
Moreover, if \(u\) is a line direction of \(f\), then so is \(ku\) for every \(k\in\F_p\setminus\{0\}\).
\end{enumerate}
\end{theorem}

\begin{corollary}\label{cor:lower}
For every odd prime \(p\) and every transversal \(T\) of \(H_p\) in \(G_p\), we have \(\abs{T-T}\ge3p^2-2\).
Hence \(\delta(G_p,H_p)\ge3p^2-2\).
\end{corollary}

\begin{corollary}\label{cor:reduction}
For each odd prime \(p\) the following are equivalent.
\begin{enumerate}
\item[(i)] \(\delta(G_p,H_p)=(2p-1)^2\), that is, Conjecture~\ref{conj:BPT} holds at \(p\).
\item[(ii)] \(\sum_{u\in\F_p^2}\abs{A_u(f)}\ge(2p-1)^2\) for every \(f\colon\F_p^2\to\F_p^2\).
\item[(iii)] \(\sum_{u\in\F_p^2}\abs{A_u(f)}\ge(2p-1)^2\) for every \(f\) with no line direction.
\item[(iv)] \(\sum_{u\ne0}\bigl(\abs{A_u(f)}-3\bigr)\ge(p-1)(p-3)\) for every \(f\) with no line direction.
\end{enumerate}
In (ii)--(iv) one may restrict to functions with \(f(0)=0\).
\end{corollary}

\begin{corollary}\label{cor:p3}
For \(p=3\),
\[
\delta\bigl((\Z/9\Z)^2,\,3(\Z/9\Z)^2\bigr)=25 .
\]
\end{corollary}

\begin{remark}\label{rem:special}
A two-element set lies on a line.
So Theorem~\ref{thm:main} applies as soon as \(\abs{A_u(f)}=2\) for a single \(u\ne0\).
It therefore contains the statement that the notes \cite{automath} call a provisional theorem (Section~\ref{sec:intro}).
A fortiori it contains \cite[Lemma~6.4]{BPT}, which assumes two linearly independent such directions.
Part (a) is also the projection bound of \cite{leshner}, obtained earlier on the same day (Section~\ref{sec:prior}).
\end{remark}

\section{Lemmas}\label{sec:lemmas}

The lemmas are proved in order.
Each proof uses only Section~\ref{sec:setting} and the lemmas before it.
Lemmas~\ref{lem:fibre} and~\ref{lem:cycle}, and the first statement of Lemma~\ref{lem:box}, are results of the source, proved again here for completeness.

\begin{lemma}[fibre formula {\cite[Lemma~6.1]{BPT}}]\label{lem:fibre}
For all \(f\) and all \(u,x\in\F_p^2\),
\begin{equation}\label{eq:fibre}
t(x+u)-t(x)=[u]+\iota\bigl(d_uf(x)\bigr)\quad\text{in }G .
\end{equation}
Consequently:
\begin{enumerate}
\item[(i)] \((T_f-T_f)\cap\pi^{-1}(u)=[u]+\iota(A_u(f))\), a set with \(\abs{A_u(f)}\) elements;
\item[(ii)] \(\abs{T_f-T_f}=\sum_{u\in\F_p^2}\abs{A_u(f)}\);
\item[(iii)] \(A_0(f)=\{0\}\).
\end{enumerate}
\end{lemma}

\begin{proof}
The proof expands both sides with the definition of the carry.
Let \(j\in\{1,2\}\) be a coordinate index, and write \(x_j\) and \(u_j\) for the \(j\)-th coordinates of \([x]\) and \([u]\).
Then \(x_j+u_j\in\{0,\dots,2p-2\}\), and \([x+u]_j\) equals \(x_j+u_j\) or \(x_j+u_j-p\).
So \([x+u]=[x]+[u]+p\tilde c\) in \(\Z^2\) with \(\tilde c\in\{0,-1\}^2\), and \(\tilde c\bmod p=c_u(x)\) by \eqref{eq:carry}.
Hence, in \(G\), using Lemma~\ref{lem:facts}(a) three times,
\begin{align*}
t(x+u)-t(x)&=[u]+p\tilde c+\iota(f(x+u))-\iota(f(x))\\
&=[u]+\iota\bigl(c_u(x)+f(x+u)-f(x)\bigr)=[u]+\iota\bigl(d_uf(x)\bigr).
\end{align*}
(i) Every element of \(T_f-T_f\) is \(t(y)-t(x)\) for some \(x,y\), and \(\pi(t(y)-t(x))=y-x\).
So the part of \(T_f-T_f\) lying in \(\pi^{-1}(u)\) is \(\{t(x+u)-t(x)\mid x\in\F_p^2\}\), which is \([u]+\iota(A_u(f))\) by \eqref{eq:fibre}.
Translation and \(\iota\) are injective, so this set has \(\abs{A_u(f)}\) elements.
(ii) The sets \(\pi^{-1}(u)\), \(u\in\F_p^2\), partition \(G\); apply (i).
(iii) \(c_0(x)=0\), so \(d_0f(x)=0\) for every \(x\).
\end{proof}

\begin{lemma}[cycle identity {\cite[Lemma~6.3]{BPT}}]\label{lem:cycle}
For all \(f\) and all \(u,x\in\F_p^2\),
\begin{equation}\label{eq:cycle}
\sum_{i=0}^{p-1}d_uf(x+iu)=-u .
\end{equation}
\end{lemma}

\begin{proof}
We evaluate a telescoping sum in \(G\) in two ways.
The sum \(\sum_{i=0}^{p-1}\bigl(t(x+(i+1)u)-t(x+iu)\bigr)\) telescopes to \(t(x+pu)-t(x)=0\), because \(pu=0\) in \(\F_p^2\).
By \eqref{eq:fibre} and Lemma~\ref{lem:facts}(a), the same sum equals \(p[u]+\iota\bigl(\sum_id_uf(x+iu)\bigr)=\iota\bigl(u+\sum_id_uf(x+iu)\bigr)\).
As \(\iota\) is injective, \(u+\sum_id_uf(x+iu)=0\).
\end{proof}

The next lemma says that a line containing \(A_u(f)\) must be parallel to \(u\).
For \(\abs{A_u(f)}\le2\) it is \cite[Lemma~6.3(1)--(2)]{BPT}; the proof is the same for a collinear set of any size.

\begin{lemma}\label{lem:parallel}
Let \(u\ne0\).
\begin{enumerate}
\item[(i)] If \(A_u(f)\subseteq a+\gen{s}\) with \(s\ne0\), then \(\gen{s}=\gen{u}\), and \(A_u(f)\subseteq a'+\gen{u}\) for every \(a'\in A_u(f)\).
\item[(ii)] \(\abs{A_u(f)}\ge2\).
\end{enumerate}
\end{lemma}

\begin{proof}
(i) We sum the points \(d_uf(iu)\) of the line and use Lemma~\ref{lem:cycle}.
Write \(d_uf(iu)=a+\lambda_is\) with \(\lambda_i\in\F_p\), for \(i=0,\dots,p-1\).
By \eqref{eq:cycle} with \(x=0\),
\[
-u=\sum_{i=0}^{p-1}(a+\lambda_is)=pa+\Lambda s=\Lambda s,\qquad\text{where } \Lambda=\sum_{i=0}^{p-1}\lambda_i ,
\]
since \(pa=0\) in \(\F_p^2\).
As \(u\ne0\), we get \(\Lambda\ne0\) and \(u=-\Lambda s\in\gen{s}\).
So \(\gen{u}\subseteq\gen{s}\).
Both sets have \(p\) elements, so they are equal.
Then \(a+\gen{s}=a+\gen{u}=a'+\gen{u}\) for each \(a'\in A_u(f)\).

(ii) The set \(A_u(f)\) is non-empty.
Suppose it were a single point \(a\).
Choose \(s\notin\gen{u}\); such \(s\) exists since \(p^2>p\).
Then \(A_u(f)\subseteq a+\gen{s}\), and (i) gives \(\gen{s}=\gen{u}\), a contradiction.
\end{proof}

\begin{lemma}\label{lem:cyclic}
Let \(g\in G\setminus H\), \(u=\pi(g)\) and \(C=\gen{g}\). Then \(u\ne0\), and:
\begin{enumerate}
\item[(i)] \(\abs{C}=p^2\);
\item[(ii)] \(\pi(C)=\gen{u}\);
\item[(iii)] \(C\cap H=pC=\iota(\gen{u})\), a subgroup of order \(p\);
\item[(iv)] \(C\cap\pi^{-1}(u)=g+pC\);
\item[(v)] every element of \(C\setminus pC\) has order \(p^2\).
\end{enumerate}
\end{lemma}

\begin{proof}
\(u\ne0\) because \(g\notin H=\ker\pi\).
(i) is Lemma~\ref{lem:facts}(c).
(ii) \(\pi(kg)=ku\) for every integer \(k\).
(iii) \(kg\in H\) if and only if \(ku=0\), if and only if \(p\) divides \(k\).
So \(C\cap H=\{jpg\mid j\in\Z\}=pC\).
By Lemma~\ref{lem:facts}(b), \(jpg=j\iota(u)=\iota(ju)\), so \(pC=\iota(\gen{u})\).
This set has \(p\) elements because \(\iota\) is injective (Lemma~\ref{lem:facts}(a)).
(iv) \(\pi(kg)=u\) if and only if \(k\equiv1\pmod p\), if and only if \((k-1)g\in pC\) by (iii), that is, \(kg\in g+pC\).
(v) By (iii), \(C\setminus pC=C\setminus H\subseteq G\setminus H\); apply Lemma~\ref{lem:facts}(c).
\end{proof}

\begin{lemma}[box]\label{lem:box}
For \(f=0\), \(\abs{T_0-T_0}=(2p-1)^2\).
Hence \(\delta(G_p,H_p)\le(2p-1)^2\) \textup{\cite[Proposition~5.3]{BPT}}.
Moreover \(\abs{A_w(0)}=2^{n(w)}\), where \(n(w)\) is the number of nonzero coordinates of \(w\in\F_p^2\).
\end{lemma}

\begin{proof}
\(T_0=\{[x]\mid x\in\F_p^2\}\) is the set of classes of \(\{0,\dots,p-1\}^2\).
So \(T_0-T_0\) is the set of classes of \(\{-(p-1),\dots,p-1\}^2\).
The \(2p-1\) integers \(-(p-1),\dots,p-1\) are pairwise incongruent modulo \(p^2\), because two of them differ by at most \(2p-2<p^2\).
So \(\abs{T_0-T_0}=(2p-1)^2\).

For the last claim, \(A_w(0)=\{c_w(x)\mid x\in\F_p^2\}\), and the \(j\)-th coordinate of \(c_w(x)\) depends only on \(x_j\) and \(w_j\).
In a coordinate \(j\) with \(w_j=0\) the carry is always \(0\).
In a coordinate \(j\) with \(w_j\ne0\) it takes exactly the two values \(0\) (for \(x_j=0\)) and \(-1\) (for \(x_j=p-1\)).
The two coordinates of \(x\) vary independently, so \(A_w(0)\) is a product of \(2-n(w)\) one-point sets and \(n(w)\) two-point sets.
\end{proof}

The last lemma is a one-dimensional lower bound for the differences of two transversals.
It is stated for a general abelian group \(\Gamma\) and a subgroup \(N\) of finite index.
The proof of Theorem~\ref{thm:main} uses it twice.

\begin{lemma}\label{lem:cross}
Let \(\Gamma\) be an abelian group, \(N\le \Gamma\) a subgroup of finite index \(q\ge1\), \(\sigma\colon \Gamma\to \Gamma/N\) the quotient map, and \(X_1,X_2\subseteq \Gamma\) two transversals of \(N\) in \(\Gamma\).
Assume
\begin{quote}
\textup{(O)} every \(e\in \Gamma\setminus N\) has order greater than \(q\) (possibly infinite).
\end{quote}
For \(r\in \Gamma/N\) put \(D_r=(X_1-X_2)\cap\sigma^{-1}(r)\). Then:
\begin{enumerate}
\item[(i)] each \(D_r\) is non-empty, and \(X_1-X_2\) is the disjoint union of the sets \(D_r\);
\item[(ii)] at most one of the sets \(D_r\) is a singleton;
\item[(iii)] \(\abs{X_1-X_2}\ge2q-1\);
\item[(iv)] if \(X_1=X_2\), then \(D_0=\{0\}\) and \(\abs{D_r}\ge2\) for every \(r\ne0\).
\end{enumerate}
\end{lemma}

\begin{proof}
The main step is part (ii): two singleton sets \(D_r\) would make \(X_1\) invariant under translation by an element of small order.
Let \(\alpha,\beta\colon \Gamma/N\to \Gamma\) be the maps with \(\sigma\circ\alpha=\sigma\circ\beta=\id\), \(\alpha(\Gamma/N)=X_1\) and \(\beta(\Gamma/N)=X_2\).
They exist and are unique because \(X_1\) and \(X_2\) are transversals.

(i) For \(x\in \Gamma/N\), \(\alpha(x+r)-\beta(x)\in D_r\), so \(D_r\ne\emptyset\).
An element \(\alpha(y)-\beta(x)\) of \(X_1-X_2\) lies in \(\sigma^{-1}(y-x)\) and in no other fibre of \(\sigma\).
So \(D_r=\{\alpha(x+r)-\beta(x)\mid x\in \Gamma/N\}\), and the sets \(D_r\) partition \(X_1-X_2\).

(ii) Suppose \(r\ne r'\), \(D_r=\{b\}\) and \(D_{r'}=\{b'\}\).
Then \(\alpha(x+r)-\beta(x)=b\) and \(\alpha(x+r')-\beta(x)=b'\) for all \(x\in \Gamma/N\).
Subtracting, \(\alpha(x+r)-\alpha(x+r')=e\) for all \(x\), where \(e=b-b'\).
Here \(\sigma(e)=r-r'\ne0\), so \(e\notin N\).
Putting \(y=x+r'\) gives \(\alpha(y+(r-r'))=\alpha(y)+e\) for all \(y\in \Gamma/N\).
So \(z+e\in X_1\) whenever \(z\in X_1\), and by induction \(z+je\in X_1\) for all \(z\in X_1\) and all integers \(j\ge0\).
Fix \(z\in X_1\).
The \(q+1\) elements \(z,z+e,\dots,z+qe\) lie in \(X_1\), which has \(q\) elements.
So \(z+je=z+j'e\) for some \(0\le j<j'\le q\).
Then \((j'-j)e=0\) with \(1\le j'-j\le q\), which contradicts~(O).

(iii) By (i) and (ii), \(\abs{X_1-X_2}=\sum_{r\in \Gamma/N}\abs{D_r}\ge1+2(q-1)=2q-1\).

(iv) If \(X_1=X_2\), then \(\alpha=\beta\) and \(D_0=\{\alpha(x)-\alpha(x)\mid x\in \Gamma/N\}=\{0\}\) is a singleton.
By (ii) no other \(D_r\) is a singleton, and by (i) none is empty.
\end{proof}

\begin{remark}\label{rem:cross}
Hypothesis (O) cannot be dropped.
In \(\Gamma=\Z/6\Z\) with \(N=\{0,3\}\) and \(q=3\), the transversal \(X_1=X_2=\{0,2,4\}\) has \(\abs{X_1-X_2}=3<5\); here the element \(2\notin N\) has order \(3=q\).
Let \(m(\Gamma,N)\) be the largest order of a subgroup of \(\Gamma\) that meets \(N\) only in \(0\), as in \cite{BPT}.
For \(\Gamma\) cyclic of order \(p^2\) and \(N=p\Gamma\) one has \(m(\Gamma,N)=1\) \cite[Proposition~5.3]{BPT}.
In that case part (iii) is a special case of the bound \(\abs{X_1-X_2}\ge2\abs{\Gamma/N}-m(\Gamma,N)\) of \cite[Theorem~4.2]{BPT}, which is proved there with Kneser's theorem.
For \(X_1=X_2\) in that case, parts (iii) and (iv) also follow from \cite[Propositions~3.2 and~3.3]{BPT}.
The proof above adapts the singleton argument of the source to two transversals and does not use Kneser's theorem.
\end{remark}

\section{\texorpdfstring{Proof of Theorem~\ref{thm:main}}{Proof of the main theorem}}\label{sec:proof}

The proof has three parts.
First, Lemma~\ref{lem:cover} turns the hypothesis into a covering of \(T\) by \(p\) cosets of a cyclic subgroup \(C\).
Second, we set up two groups of order \(p^2\) with subgroups of index \(p\) in which Lemma~\ref{lem:cross} applies.
Third, we apply Lemma~\ref{lem:cross} in both and count the differences coset by coset.

\subsection{The covering criterion}

\begin{lemma}\label{lem:cover}
Let \(u\ne0\), \(a\in\F_p^2\), \(g=[u]+\iota(a)\) and \(C=\gen{g}\).
Then \(C\) is cyclic of order \(p^2\) with \(\pi(C)=\gen{u}\), and for every \(f\) the following are equivalent:
\begin{enumerate}
\item[(1)] \(A_u(f)\subseteq a+\gen{u}\);
\item[(2)] \(T_f\) is contained in a union of at most \(p\) cosets of \(C\).
\end{enumerate}
When they hold, \(t\) maps each \(u\)-line into one coset of \(C\), and different \(u\)-lines go to different cosets.
So \(T_f\) meets exactly \(p\) cosets of \(C\), each in exactly \(p\) elements.
Moreover, let \(C'\le G\) be any cyclic subgroup of order \(p^2\).
Then for each nonzero \(w\in\pi(C')\) there is \(a'\in\F_p^2\) with \(C'=\gen{[w]+\iota(a')}\).
\end{lemma}

\begin{proof}
Since \(\pi(g)=u\ne0\), we have \(g\notin H\), and Lemma~\ref{lem:cyclic} applies to \(g\) and \(C\).
It gives the first sentence.
For the equivalence we show that (1) puts every step \(t(x+u)-t(x)\) into \(C\), and that (2) forces the converse.

\emph{Step 1 (one step).}
Assume (1) and let \(x\in\F_p^2\).
We use \eqref{eq:fibre}, Lemma~\ref{lem:facts}(a) and Lemma~\ref{lem:cyclic}(iii):
\begin{align*}
t(x+u)-t(x)&=[u]+\iota(d_uf(x))=g+\iota\bigl(d_uf(x)-a\bigr)\\
&\in g+\iota(\gen{u})=g+pC\subseteq C .
\end{align*}

\emph{Step 2 (chains).}
For \(x\in\F_p^2\) and an integer \(i\ge0\), the difference \(t(x+iu)-t(x)\) is the sum of the \(i\) steps \(t(x+(j+1)u)-t(x+ju)\), \(0\le j<i\).
So it lies in \(C\) by Step~1.
Hence \(t\) maps the \(u\)-line \(x+\gen{u}\) into the coset \(t(x)+C\).
There are \(p\) such lines, so \(T_f\) lies in at most \(p\) cosets of \(C\).
This proves that (1) implies (2).

\emph{Step 3 (converse).}
Assume (2).
If \(t(x)\) and \(t(x')\) lie in the same coset of \(C\), then \(x-x'=\pi(t(x)-t(x'))\in\pi(C)=\gen{u}\).
So \(x\) and \(x'\) lie on the same \(u\)-line.
Hence the elements of \(T_f\) in one coset of \(C\) are images of points of a single \(u\)-line, and there are at most \(p\) of them.
As \(\abs{T_f}=p^2\) and at most \(p\) cosets are met, exactly \(p\) cosets are met.
Each of them contains exactly \(p\) elements of \(T_f\), namely the images of all points of one \(u\)-line.
Different cosets correspond to different \(u\)-lines.
In particular, for every \(x\) the points \(x\) and \(x+u\) are on one \(u\)-line, so \(t(x+u)-t(x)\in C\).
Its image under \(\pi\) is \(u\).
So by Lemma~\ref{lem:cyclic}(iv), Lemma~\ref{lem:cyclic}(iii) and the definition of \(g\),
\[
[u]+\iota(d_uf(x))=t(x+u)-t(x)\in g+pC=[u]+\iota\bigl(a+\gen{u}\bigr).
\]
Since \(\iota\) is injective, \(d_uf(x)\in a+\gen{u}\) for all \(x\).
This proves that (2) implies (1).
Under (1) we have (2), so the counting in this step applies; it proves the statement about exactly \(p\) cosets.

\emph{Step 4 (the last statement).}
Let \(g_0\) be a generator of \(C'\).
Then \(g_0\notin H\), since otherwise \(pg_0=0\) by Lemma~\ref{lem:facts}(b), and \(g_0\) would have order at most \(p\).
Put \(u_0=\pi(g_0)\ne0\), so \(\pi(C')=\gen{u_0}\).
Let \(w=ku_0\) with \(k\in\{1,\dots,p-1\}\), and put \(g'=kg_0\in C'\).
Then \(\pi(g')=w\ne0\), so \(g'\notin H\), and \(g'\) has order \(p^2\) by Lemma~\ref{lem:facts}(c).
Hence \(g'\) generates \(C'\).
As \(\pi(g'-[w])=0\), we have \(g'-[w]\in H=\iota(\F_p^2)\), so \(g'=[w]+\iota(a')\) for some \(a'\in\F_p^2\).
\end{proof}

\subsection{Proof of part (a)}
Let \(u\ne0\) be a line direction of \(f\).
Choose \(a\in A_u(f)\).
By Lemma~\ref{lem:parallel}(i), \(A_u(f)\subseteq a+\gen{u}\).
Let \(g=[u]+\iota(a)\) and \(C=\gen{g}\), as in Lemma~\ref{lem:cover}; condition (1) of that lemma holds.
Fix \(v_0\in\F_p^2\setminus\gen{u}\).
Since \(u,v_0\) is a basis of \(\F_p^2\), the \(u\)-lines are \(L_n=nv_0+\gen{u}\) for \(n\in\F_p\).
Put \(\tau_n=t(nv_0)\) and \(R_n=t(L_n)\).
By Lemma~\ref{lem:cover}, \(t\) maps \(L_n\) into one coset of \(C\), which contains \(\tau_n\).
Hence
\[
R_n\subseteq\tau_n+C,\qquad T=\bigsqcup_{n\in\F_p}R_n,\qquad \abs{R_n}=p .
\]
The union is disjoint and \(\abs{R_n}=p\) because \(t\) is injective (as \(\pi\circ t\) is the identity) and the lines \(L_n\) partition \(\F_p^2\).

We now define the two groups to which Lemma~\ref{lem:cross} will be applied.
Let \(M=H+C\).
Let \(\rho\colon G\to Q=G/C\) be the quotient map, \(K=\rho(H)\le Q\), and \(\kappa\colon Q\to Q/K\) the quotient map.
We prove four claims about these groups.

\emph{Claim A.} \(M=\pi^{-1}(\gen{u})\), and \(\kappa\circ\rho\colon G\to Q/K\) is a surjective homomorphism with kernel \(M\).
Hence there is an isomorphism \(\theta\colon Q/K\to\F_p^2/\gen{u}\) with \(\theta(\kappa\rho(y))=\pi(y)+\gen{u}\) for all \(y\in G\), and \(\abs{Q/K}=p\).

\begin{proof}[Proof of Claim A]
\(\pi(M)=\pi(H)+\pi(C)=\gen{u}\) by Lemma~\ref{lem:cyclic}(ii).
Conversely, if \(\pi(y)=ku\) for an integer \(k\), then \(\pi(y-kg)=0\), so \(y-kg\in H\) and \(y\in H+C\).
Next, \(\kappa\rho(y)=0\) if and only if \(\rho(y)\in\rho(H)\), if and only if \(y\in H+C=M\).
Both \(\kappa\rho\) and \(y\mapsto\pi(y)+\gen{u}\) are surjective homomorphisms from \(G\) with kernel \(M\), which gives \(\theta\).
Finally, \(\F_p^2/\gen{u}\) has \(p\) elements.
\end{proof}

\emph{Claim B.} Hypothesis (O) of Lemma~\ref{lem:cross} holds for \((\Gamma,N)=(Q,K)\) with \(q=p\).

\begin{proof}[Proof of Claim B]
Let \(e\in Q\setminus K\) and write \(e=\rho(y)\) with \(y\in G\).
Then \(y\notin M\): otherwise \(y=h+\gamma\) with \(h\in H\) and \(\gamma\in C\), and \(e=\rho(y)=\rho(h)\in K\).
Suppose \(je=0\), that is, \(jy\in C\), for some integer \(1\le j\le p\).
If \(j<p\), then \(j\pi(y)=\pi(jy)\in\pi(C)=\gen{u}\).
As \(j\) is invertible in \(\F_p\), this gives \(\pi(y)\in\gen{u}\), so \(y\in M\) by Claim~A, a contradiction.
If \(j=p\), then \(py\in H\) because \(\pi(py)=p\pi(y)=0\).
So \(py\in C\cap H=pC\) by Lemma~\ref{lem:cyclic}(iii).
Write \(py=p\gamma\) with \(\gamma\in C\).
Then \(p(y-\gamma)=0\), so \(y-\gamma\in H\) by Lemma~\ref{lem:facts}(b), and \(y\in H+C=M\), a contradiction.
\end{proof}

\emph{Claim C.} Put \(y_n=\rho(\tau_n)\) and \(Y=\{y_n\mid n\in\F_p\}\). Then \(\rho(R_n)=\{y_n\}\), and \(Y\) is a transversal of \(K\) in \(Q\).

\begin{proof}[Proof of Claim C]
\(R_n\subseteq\tau_n+C\) gives \(\rho(R_n)=\{y_n\}\).
By Claim~A, \(\theta(\kappa(y_n))=\pi(\tau_n)+\gen{u}=nv_0+\gen{u}\).
These \(p\) cosets are pairwise different: if \((n-n')v_0\in\gen{u}\), then \(n=n'\), because \(v_0\notin\gen{u}\).
As \(Q/K\) has \(p\) elements (Claim~A), \(\kappa\) maps \(Y\) bijectively onto \(Q/K\).
\end{proof}

\emph{Claim D.} \(\abs{C/pC}=p\); for each \(n\) the set \(S_n=R_n-\tau_n\) is a transversal of \(pC\) in \(C\); and hypothesis (O) holds for \((\Gamma,N)=(C,pC)\) with \(q=p\).

\begin{proof}[Proof of Claim D]
The restriction of \(\pi\) to \(C\) is a homomorphism onto \(\gen{u}\) with kernel \(C\cap H=pC\) (Lemma~\ref{lem:cyclic}(ii),~(iii)).
So \(\abs{C/pC}=p\), and two elements of \(C\) lie in the same coset of \(pC\) if and only if they have the same image under \(\pi\).
Since \(R_n\subseteq\tau_n+C\), we have \(S_n=\{t(nv_0+iu)-t(nv_0)\mid i\in\F_p\}\subseteq C\).
The element with index \(i\) has image \(iu\) under \(\pi\).
These images are pairwise different and exhaust \(\gen{u}\), so \(S_n\) has exactly one element in each coset of \(pC\).
Hypothesis (O) is Lemma~\ref{lem:cyclic}(v), since \(p^2>p\).
\end{proof}

\emph{Applying Lemma~\ref{lem:cross}.}
By Claims~A, B and~C, Lemma~\ref{lem:cross}(iii) and~(iv) apply to \((\Gamma,N,X_1,X_2)=(Q,K,Y,Y)\) with \(q=p\).
They give
\begin{equation}\label{eq:YY}
\abs{Y-Y}\ge2p-1,\qquad \abs{(Y-Y)\cap\kappa^{-1}(r)}\ge2\quad\text{for every } r\in Q/K\setminus\{0\}.
\end{equation}
By Claim~D, Lemma~\ref{lem:cross}(iii) applies to \((\Gamma,N,X_1,X_2)=(C,pC,S_n,S_m)\) with \(q=p\).
It gives
\begin{equation}\label{eq:SS}
\abs{S_n-S_m}\ge2p-1\qquad\text{for all } n,m\in\F_p .
\end{equation}

\emph{Counting.}
For \(n,m\in\F_p\) we have \(R_n=\tau_n+S_n\) and \(S_n,S_m\subseteq C\).
Hence, using \eqref{eq:SS},
\begin{align*}
R_n-R_m&=(\tau_n-\tau_m)+(S_n-S_m)\subseteq(T-T)\cap\rho^{-1}(y_n-y_m),\\
\abs{R_n-R_m}&=\abs{S_n-S_m}\ge2p-1 .
\end{align*}
Every \(z\in Y-Y\) equals \(y_n-y_m\) for some \(n,m\).
So
\begin{equation}\label{eq:fibrecount}
\abs{(T-T)\cap\rho^{-1}(z)}\ge2p-1\qquad\text{for every } z\in Y-Y .
\end{equation}
The sets \(\rho^{-1}(z)\) are pairwise disjoint, and \eqref{eq:YY} gives
\begin{align}
\abs{T-T}&\ge\sum_{z\in Y-Y}\abs{(T-T)\cap\rho^{-1}(z)}\notag\\
&\ge(2p-1)\abs{Y-Y}\ge(2p-1)^2 .\label{eq:count}
\end{align}
Together with Lemma~\ref{lem:fibre}(ii), this is part (a).
\qed

\subsection{Proof of part (b)}
We count the elements of \(T-T\) whose image under \(\pi\) lies in one coset of \(\gen{u}\), using \eqref{eq:YY} and \eqref{eq:fibrecount}.
Keep the notation of the proof of part (a).
Let \(W\) be a coset of \(\gen{u}\) in \(\F_p^2\).
By Lemma~\ref{lem:fibre}(i), \(\sum_{w\in W}\abs{A_w(f)}=\abs{(T-T)\cap\pi^{-1}(W)}\).
By Claim~A, \(\pi^{-1}(W)=(\kappa\rho)^{-1}(r)\) with \(r=\theta^{-1}(W)\), and \(r\ne0\) exactly when \(W\ne\gen{u}\).

First let \(W=v+\gen{u}\) with \(v\notin\gen{u}\).
Then \(r\ne0\), and by \eqref{eq:YY} the set \((Y-Y)\cap\kappa^{-1}(r)\) contains two different elements \(z\) and \(z'\).
The sets \((T-T)\cap\rho^{-1}(z)\) and \((T-T)\cap\rho^{-1}(z')\) are disjoint.
Each has at least \(2p-1\) elements by \eqref{eq:fibrecount}, and both lie in \((\kappa\rho)^{-1}(r)=\pi^{-1}(W)\).
Hence \(\sum_{k\in\F_p}\abs{A_{ku+v}(f)}\ge2(2p-1)=4p-2\).

Now let \(W=\gen{u}\).
Then \(R_0-R_0\subseteq(T-T)\cap\rho^{-1}(0)\subseteq(T-T)\cap M\), and \(R_0-R_0\) has at least \(2p-1\) elements by the counting step.
As \(M=\pi^{-1}(\gen{u})\) by Claim~A, \(\sum_{k\in\F_p}\abs{A_{ku}(f)}\ge2p-1\).
\qed

\begin{remark}
Summing the bounds of part (b) over the \(p\) cosets of \(\gen{u}\) gives \((2p-1)+(p-1)(4p-2)=(2p-1)^2\), the bound of part (a).
So the bounds of part (b) imply the bound of part (a).
\end{remark}

\subsection{Proof of parts (c), (d) and (e)}
Part (c) is the implication from (1) to (2) in Lemma~\ref{lem:cover}, with the subgroup \(C\) used in the proof of part (a); that lemma also gives \(\pi(C)=\gen{u}\).

For part (d), let \(T\) lie in a union of at most \(p\) cosets of a cyclic subgroup \(C'\) of order \(p^2\), and let \(w\in\pi(C')\) be nonzero.
By the last statement of Lemma~\ref{lem:cover}, \(C'=\gen{[w]+\iota(a')}\) for some \(a'\in\F_p^2\).
The implication from (2) to (1) in Lemma~\ref{lem:cover}, with \(u,a\) replaced by \(w,a'\), gives \(A_w(f)\subseteq a'+\gen{w}\), which is an affine line.

For part (e), first let \(f\) have a line direction \(u\).
By part (c), \(T\) lies in a union of \(p\) cosets of a cyclic subgroup of order \(p^2\).
Conversely, let \(T\) lie in a union of at most \(p\) cosets of a cyclic subgroup \(C'\) of order \(p^2\).
Then \(\pi(C')\ne\{0\}\).
Indeed, otherwise \(C'\subseteq\ker\pi=H\); a generator \(g'\) of \(C'\) would satisfy \(pg'=0\) by Lemma~\ref{lem:facts}(b), and \(C'\) would have at most \(p\) elements.
So \(\pi(C')\) has a nonzero element, and it is a line direction of \(f\) by part (d).
Finally, let \(u\) be a line direction.
By part (c), \(T\) lies in \(p\) cosets of a cyclic subgroup \(C\) of order \(p^2\) with \(\pi(C)=\gen{u}\).
By part (d), every nonzero \(ku\in\gen{u}\) is a line direction.
\qed

\subsection{Remarks}

\begin{remark}[sharpness]\label{rem:sharp}
Take the box \(f=0\) and \(u=e_1=(1,0)\).
By Lemma~\ref{lem:box} and its proof, \(A_{e_1}(0)=\{(0,0),(-1,0)\}\subseteq\gen{e_1}\), and \(\abs{T_0-T_0}=(2p-1)^2\).
For \(v=(0,\lambda)\) with \(\lambda\ne0\), Lemma~\ref{lem:box} gives the column sum \(2+4(p-1)=4p-2\) for \(v+\gen{e_1}\).
So the bounds of parts (a) and (b) are attained for every \(p\).
\end{remark}

\begin{remark}[the prime \(2\)]\label{rem:p2}
The proofs of the lemmas and of Theorem~\ref{thm:main} never use that \(p\) is odd.
At \(p=2\), Lemma~\ref{lem:cycle} gives \(d_uf(x+u)=d_uf(x)+u\).
So each \(A_u(f)\) with \(u\ne0\) is a union of cosets of \(\gen{u}\) and has \(2\) or \(4\) elements.
If some \(A_u(f)\) has \(2\) elements, it is a line and Theorem~\ref{thm:main}(a) gives \(\abs{T-T}\ge9\); otherwise \(\abs{T-T}=1+3\cdot4=13\).
So every transversal of \(2(\Z/4\Z)^2\) in \((\Z/4\Z)^2\) has at least \(9\) differences.
This agrees with \cite[Corollary~5.2]{BPT} and serves only as a consistency check.
\end{remark}

\section{Proofs of the corollaries}\label{sec:corollaries}

\begin{proof}[Proof of Corollary~\ref{cor:lower}]
We split according to whether \(f\) has a line direction.
Let \(T\) be a transversal; then \(T=T_f\) for some \(f\) by Lemma~\ref{lem:Tf}.
If \(f\) has a line direction, then \(\abs{T-T}\ge(2p-1)^2\) by Theorem~\ref{thm:main}(a).
Otherwise no \(A_u(f)\) with \(u\ne0\) lies on an affine line.
Every set of at most two points lies on an affine line (Section~\ref{sec:setting}), so \(\abs{A_u(f)}\ge3\) for all \(u\ne0\).
By Lemma~\ref{lem:fibre}(ii) and~(iii),
\[
\abs{T-T}=1+\sum_{u\ne0}\abs{A_u(f)}\ge1+3(p^2-1)=3p^2-2 .
\]
In the first case, \(\abs{T-T}\ge(2p-1)^2\ge3p^2-2\), because \((2p-1)^2-(3p^2-2)=p^2-4p+3=(p-1)(p-3)\ge0\) for \(p\ge3\).
So \(\abs{T-T}\ge3p^2-2\) in both cases.
\end{proof}

\begin{remark}\label{rem:compare}
For every odd prime \(p\), the bound \(3p^2-2\) exceeds the bound \(3p^2-p-1\) of \cite[Theorem~6.5]{BPT} by \(p-1\).
At \(p=3\) the two bounds are \(25\) and \(23\), and the two lower bounds \((2p-1)^2\) and \(3p^2-2\) in the proof above are both equal to \(25\).
\end{remark}

\begin{proof}[Proof of Corollary~\ref{cor:p3}]
Let \(p=3\).
Then Corollary~\ref{cor:lower} gives \(\delta(G_3,H_3)\ge3\cdot3^2-2=25\).
Lemma~\ref{lem:box} gives \(\delta(G_3,H_3)\le25\).
This proof uses no computation.
\end{proof}

\begin{proof}[Proof of Corollary~\ref{cor:reduction}]
(i)\(\Leftrightarrow\)(ii): the transversals are exactly the sets \(T_f\) (Lemma~\ref{lem:Tf}), and \(\abs{T_f-T_f}=\sum_u\abs{A_u(f)}\) (Lemma~\ref{lem:fibre}(ii)).
So (ii) says \(\delta(G_p,H_p)\ge(2p-1)^2\), and the reverse inequality always holds (Lemma~\ref{lem:box}).

(ii)\(\Rightarrow\)(iii) is immediate.

(iii)\(\Rightarrow\)(ii): if \(f\) has a line direction, the inequality is Theorem~\ref{thm:main}(a); otherwise it is (iii).

(iii)\(\Leftrightarrow\)(iv): by Lemma~\ref{lem:fibre}(iii),
\[
\sum_u\abs{A_u(f)}=1+3(p^2-1)+\sum_{u\ne0}\bigl(\abs{A_u(f)}-3\bigr),
\]
and \((2p-1)^2-1-3(p^2-1)=(p-1)(p-3)\).

Normalisation: for a constant \(b\in\F_p^2\), \(d_u(f-b)=d_uf\) for every \(u\).
So \(f\) and \(f-f(0)\) have the same sets \(A_u\) and the same line directions.
\end{proof}

\section{Computer checks and formalization}\label{sec:computer}

No step of the proof above uses a computation, and no argument depends on the checks described here.
As evidence only, an exact-arithmetic Python script tests individual steps directly from the definitions.
It recomputes \(T_f-T_f\) by literal subtraction in \(G\), computes \(d_uf\) from \eqref{eq:carry} and~\eqref{eq:du}, and builds subgroups and cosets as explicit sets.
At \(p=3\) it covers all 704{,}214 functions with \(f(0)=0\) and a line direction; among them the least value of \(\abs{T-T}\) is 25 and the least column sum is 10.
A separate exhaustive scan of all \(9^8\) functions with \(f(0)=0\) finds the same number of such functions.
The script also tests the statements about \(G\) alone, for every element and every cyclic subgroup of order \(p^2\), at \(p=2,3,5,7,11\); Lemma~\ref{lem:cross} on all pairs of transversals in \(\Z/n\Z\) for \(n=1,4,8,9,16,25,27\); all 64 functions with \(f(0)=0\) at \(p=2\); and stored and random functions at \(p=3,5,7,11\).
Every archived run ends with the line \texttt{RESULT: ALL CHECKS PASS}.
The scripts, the logs, the exact commands and a table of which steps are tested are in the accompanying repository \cite{repo}.
Nothing in this paper has been formalized in a proof assistant.

\section*{Use of AI}

Human directed research produced with the help of AI tools (Claude) and checked by AI referees (Claude Opus and Claude Fable; ChatGPT and Gemini as outside cross-checks).
However, additional verification may still be required.
Code and full proof: \cite{repo}.

\begin{thebibliography}{9}\raggedright

\bibitem{BPT}
M.~Barcau, V.~Pa\c{s}ol, and G.~C.~\c{T}urca\c{s}, \emph{Transversal difference numbers in finite abelian quotients}, preprint, arXiv:2606.27961v1, 2026.
\url{https://arxiv.org/abs/2606.27961}.

\bibitem{BPTcode}
Companion repository of \cite{BPT}, \emph{georgeturcasubb/transversals}, 2026; read on 2026-11-27.
\url{https://github.com/georgeturcasubb/transversals}.

\bibitem{leshner}
R.~Leshner, with Codex, \emph{Square-plane transversal difference number at \(p=7\) is \(169\)}, GitHub pull request~\#6492 to GuillaumeBAUD-coder/open-math-agents, opened 2026-11-27 (12:56 UTC); commit \texttt{ba804a2}, read on 2026-11-27.
\url{https://github.com/GuillaumeBAUD-coder/open-math-agents/pull/6492};
README, section ``A projection bound for every odd prime'':
\url{https://github.com/rleshner/open-math-agents/blob/ba804a2813650e214c5172a1f484b1c1eda5edd8/research/transversal-seven/README.md}.

\bibitem{automath}
chy4pro/automath, problem folder \texttt{tdn\_2606\_27961}, unpublished working notes, 2026; commit \texttt{5542184}, read on 2026-11-27.
\url{https://github.com/chy4pro/automath/tree/5542184180d014b68637851941109dd5e584e920/problems/tdn_2606_27961}.
Status note, lines 36--40:
\url{https://github.com/chy4pro/automath/blob/5542184180d014b68637851941109dd5e584e920/notes/selection/selection_0824.md}.

\bibitem{trureturing}
the-omega-institute/trureturing, \emph{Research: square-plane transversal difference Conjecture~6.2}, issue~\#8487, opened 2026-09-17; read on 2026-11-27.
\url{https://github.com/the-omega-institute/trureturing/issues/8487}.

\bibitem{repo}
S.~Sachdeva, \emph{Accompanying repository of this paper: code, archived logs, full proof and review history}, version of 2026-11-27.
\url{https://github.com/nebulousice/barcau-pasol-turcas-conjecture-6-2}.

\end{thebibliography}
\end{document}
